\documentclass[11pt]{article}

\usepackage[T1]{fontenc}
\usepackage[utf8]{inputenc}
\usepackage{lmodern}
\usepackage[margin=1in]{geometry}
\usepackage{amsmath,amssymb,amsthm,mathtools}
\usepackage{booktabs,tabularx}
\usepackage[shortlabels]{enumitem}
\usepackage{xcolor}
\usepackage{hyperref}
\usepackage{xurl}
\usepackage{microtype}

\hypersetup{colorlinks=true,linkcolor=blue,citecolor=blue,urlcolor=blue}
\setlength{\parindent}{0pt}
\setlength{\parskip}{0.55em}
\setlength{\emergencystretch}{3em}
\numberwithin{equation}{section}

\newtheorem{theorem}{Theorem}[section]
\newtheorem{proposition}[theorem]{Proposition}
\newtheorem{corollary}[theorem]{Corollary}
\newtheorem{assumption}[theorem]{Assumption}
\theoremstyle{remark}
\newtheorem{remark}[theorem]{Remark}

\newcommand{\OPH}{\mathrm{OPH}}
\newcommand{\RPC}{\mathrm{RPC}}
\newcommand{\chinu}{\chi_\nu}
\newcommand{\chican}{\chi_\nu^{\mathrm{can}}}
\newcommand{\chieng}{\chi_\nu^{\mathrm{eng}}}
\newcommand{\scohcan}{S_{\mathrm{coh}}^{\mathrm{can}}}
\newcommand{\scoheng}{S_{\mathrm{coh}}^{\mathrm{eng}}}
\newcommand{\Ncoh}{N_{\mathrm{coh}}}
\newcommand{\lambdac}{\lambda_{\mathrm{collar}}}
\newcommand{\dd}{\,\mathrm d}

\title{\textbf{Theoretical Bounds on \(\chinu\) in Observer-Patch Holography}\\
\large Canonical Coherent-Matter Susceptibility and Repair-Charge Coupling}
\author{Bernhard Mueller, Alex Osika, and David Matscheko}
\date{\OPHPaperReleaseDate}
\newcommand{\OPHPaperReleaseID}{r1520}
\newcommand{\OPHPaperReleaseDate}{July 9, 2026}
\newcommand{\OPHPaperReleaseBanner}{%
\begin{center}
\small\textbf{Paper release:} \texttt{\OPHPaperReleaseID}\quad\textbf{Released:} \OPHPaperReleaseDate
\end{center}
}


\begin{document}
\maketitle
\begin{center}
\small\textbf{Paper release:} \OPHPaperReleaseID\quad
\textbf{Released:} \OPHPaperReleaseDate
\end{center}

\begin{abstract}
The coherent-material branch of observer-patch holography supplies a canonical
scalar source in the same quotient-edge register as the protected collar
reserve. On the co-registered presence branch, the dimensionless coupling is
\(\chi_\nu^{\rm can}=1-P/24=0.9320429912748350\ldots\). Under the weaker
mean-count reading it lies between that value and one. This is a channel
coefficient, not a force coefficient. A proposed repair-charge condensate
action couples the source to a dynamical repair phase. Conditional on that
action, repair-current continuity, dilute dust behavior, spherical deep-galaxy
scaling, field stress, and a lawful coherent-source force follow from one
variational system. A compact neutral coherence contrast has no monopole and
cannot support a closed device. The dimensional source calibration, ambient
repair field, compact-phase completion, and laboratory source receipt are work
in progress; no numerical lift prediction follows from \(\chi_\nu\) alone.
\end{abstract}

\section*{Claim Boundary}

This paper proves a finite-channel coefficient on a declared OPH branch. It
also derives continuity, dilute dust behavior, spherical deep-galaxy scaling,
force, and momentum formulae conditional on the repair-charge condensate
action in the dark-sector paper~\cite{oph-dark}. The OPH derivation of that
action and its dimensional source normalization are work in progress.

The exact statuses are:
\begin{center}
\begin{tabularx}{0.94\textwidth}{@{}lX@{}}
\toprule
Object & Status \\
\midrule
Coherent source in the quotient-edge scalar register & finite theorem under the stated source receipts \\
\(\chican=\lambdac\) & theorem under co-registration and reserve-survival premises \\
\(\chican=1-P/24\) & exact on the presence branch \\
Repair-charge action & proposed dynamical completion \\
Continuity, dust limit, and spherical deep-galaxy scaling & derived conditional on that action and its constitutive premises \\
Dimensional coupling \(q_\star\) & work in progress \\
Device force & conditional on source charge, external field, and momentum closure \\
Numerical lift & absent \\
\bottomrule
\end{tabularx}
\end{center}

\tableofcontents

\section{Canonical Coherent-Matter Source}

An OPH technology instantiation is an observer-like self-reading system: a
bounded physical or software patch with local state, ports or boundaries,
readback, records, feedback or repair moves, and a public evidence bundle. For
a nondegenerate coherent material subfederation \(U\), define the canonical
source
\begin{equation}
\scohcan
=\mathbf 1_{\rm self\mbox{-}read}\,\mathsf R_U\mathsf P_U\mathsf C_U.
\label{eq:source}
\end{equation}
The factors state that the patch reads itself, produces a stable record,
supports a predictive boundary, and executes a registered repair or control
move. If any factor vanishes, the coherent scalar source is zero.

The screen-microphysics branch supplies four structural facts
\cite{oph-micro}:
\begin{enumerate}[label=S\arabic*.,leftmargin=*]
\item \(\scohcan\) descends to a protected quotient observable.
\item A nondegenerate source is nonzero.
\item Scalar edge-center exhaustion places it in the unique quotient-local
scalar register.
\item The protected \(\mathbb Z_6\) reserve acts on that same register and
commutes with scalar activation.
\end{enumerate}
These statements identify the channel. They supply no spacetime action or
dimensional gravitational charge.

\section{Protected-Reserve Coefficient}

Let \(\epsilon_{\mathbb Z_6}(y)\) be the conditional reserve-presence profile
on the co-registered scalar slot, and let \(w(y)\) be the normalized scalar
weight across the finite collar. The scalar opportunity survives with local
factor
\begin{equation}
\lambda_{\rm slot}(y)=1-\epsilon_{\mathbb Z_6}(y).
\end{equation}
The finite-thickness collar coefficient is
\begin{equation}
\lambdac=\int\dd y\,w(y)
\left[1-\epsilon_{\mathbb Z_6}(y)\right].
\label{eq:collar}
\end{equation}
The protected-reserve receipt gives
\begin{equation}
\int\dd y\,w(y)\epsilon_{\mathbb Z_6}(y)=\frac{P}{24}.
\label{eq:reserve}
\end{equation}

\begin{theorem}[Canonical susceptibility]
Assume the nonzero coherent-source branch, scalar edge-center exhaustion,
co-registration with the protected reserve, and the presence reading in
Eqs.~\eqref{eq:collar}--\eqref{eq:reserve}. Then
\begin{equation}
\boxed{\chican=\lambdac=1-\frac{P}{24}}.
\label{eq:canonical}
\end{equation}
For \(P=1.630968209403959\),
\begin{equation}
\boxed{\chican=0.9320429912748350\ldots}.
\label{eq:value}
\end{equation}
\end{theorem}

\begin{proof}
The coherent source enters the scalar opportunity count before the protected
reserve acts. Co-registration makes the local surviving fraction
\(1-\epsilon_{\mathbb Z_6}(y)\). Averaging with \(w(y)\) gives
\(\lambdac\scohcan\). The definition of canonical susceptibility gives
\(\chican\scohcan\). Nondegeneracy permits cancellation of \(\scohcan\), and
Eq.~\eqref{eq:reserve} gives the value.
\end{proof}

\begin{corollary}[Mean-count band]
If \(\epsilon_{\mathbb Z_6}(y)\) is interpreted only as the mean of a
nonnegative integer occupancy, then
\begin{equation}
\boxed{0.9320429912748350\ldots\le\chican\le1}.
\label{eq:band}
\end{equation}
\end{corollary}

\begin{proof}
Markov's inequality gives
\(1-\epsilon\le\Pr[N=0]\le1\). Integration and
Eq.~\eqref{eq:reserve} give the band.
\end{proof}

The theorem excludes a zero scalar-channel coefficient on the stated
nondegenerate branch. It does not exclude a zero dimensional gravitational
coupling \(q_\star\), a zero device repair charge, or a zero ambient repair
gradient.

\section{Engineering Coordinates}

An engineering description may absorb an effective coherent capacity into
the source coordinate:
\begin{equation}
\scoheng=\Ncoh\scohcan.
\end{equation}
Invariance of the scalar perturbation requires
\begin{equation}
\chican\scohcan=\chieng\scoheng,
\qquad
\boxed{\chieng=\frac{\chican}{\Ncoh}}.
\label{eq:chart}
\end{equation}
Very small values of \(\chieng\) can therefore represent an order-one
canonical coefficient. Equation~\eqref{eq:chart} is a chart identity. It does
not define \(q_\star\), measure \(\scohcan\), or predict force.

\section{Repair-Charge Coupling}

The dark-sector completion promotes integer scalar repair occupation and a
compact repair phase to a canonical pair \((n,\theta)\). On a fixed real lift
of the phase, its relevant action is
\begin{align}
I_{\RPC}={}&I_{\rm Regge}+I_{\rm SM}
\nonumber\\
&+\int\dd\tau\left[
-\sum_vV_vn_vD_\tau\theta_v
-\sum_vV_v\varepsilon(n_v)
-\sum_eW_e\mathcal K(|(d\theta)_e|)
\right.\nonumber\\
&\hspace{31mm}\left.
-\sum_vV_v\theta_v
\left(\mathcal Q_{b,v}
+q_\star\chican S_{{\rm coh},v}^{\rm can}\right)
\right].
\label{eq:rpc-action}
\end{align}
The source and canonical signs are chosen so phase variation gives the
positive-source balance in Eq.~\eqref{eq:continuity}, and translation of a
positive localized source gives force \(-Q_R\nabla\theta\). A globally
single-valued compact-phase coupling, the derivation of this action from finite
repair dynamics, and the dimensional constants remain open.

The coherent repair-charge density is
\begin{equation}
\mathcal Q_{\rm coh}=q_\star\chican\scohcan.
\label{eq:repair-charge-density}
\end{equation}
The phase equation is
\begin{equation}
D_\tau n+d^\dagger j_R
=\mathcal Q_b+q_\star\chican\scohcan.
\label{eq:continuity}
\end{equation}
For a finite region \(\Omega\),
\begin{equation}
\frac{\dd}{\dd\tau}\sum_{v\in\Omega}V_vn_v
+\sum_{e\in\partial\Omega}W_ej_{R,e}
=\sum_{v\in\Omega}V_v
\left(\mathcal Q_{b,v}+q_\star\chican
S_{{\rm coh},v}^{\rm can}\right).
\label{eq:finite-balance}
\end{equation}
The boundary flux is part of the physical source ledger.

\begin{proposition}[Conditional rotor consequences]
\label{prop:rotor-consequences}
Assume Eq.~\eqref{eq:rpc-action}, a positive branch \(n>0\), and the infrared
constitutive functions
\begin{equation}
\varepsilon(n)=m_Rn+\frac{u_R}{3}n^3,
\qquad
\mathcal K(y)=\frac{\kappa_R}{3}y^3,
\qquad
m_R,u_R,\kappa_R>0.
\label{eq:rpc-constitutive}
\end{equation}
Then:
\begin{enumerate}[label=(\roman*),leftmargin=*]
\item phase variation gives the exact finite repair-current balance
Eq.~\eqref{eq:finite-balance};
\item the homogeneous dilute phase has
\begin{equation}
\rho_R=m_Rn+\frac{u_R}{3}n^3,
\qquad
p_R=\frac{2u_R}{3}n^3,
\end{equation}
so source-free expansion gives \(n\propto a^{-3}\),
\(\rho_R\propto a^{-3}\), and \(w_R\simeq0\) when
\(u_Rn^2\ll m_R\);
\item if \(\mathcal Q_b=\beta_b\rho_b\) and
\(\mathbf a_R=-\beta_b\nabla\theta\), the static spherical exterior branch
gives
\begin{equation}
a_R=\sqrt{a_ba_0},
\qquad
v^4=GM_ba_0,
\qquad
a_0=\frac{\beta_b^3}{4\pi G\kappa_R};
\end{equation}
\item translation of the coherent source gives
Eq.~\eqref{eq:force} and the momentum ledger
Eq.~\eqref{eq:momentum}.
\end{enumerate}
These statements are consequences of the proposed action, not derivations of
that action from the finite OPH theorems.
\end{proposition}

\begin{proof}
Item (i) follows by varying \(\theta\) and summing the local equation over a
finite region. For item (ii), the first-order rotor Hamiltonian density is
\(\varepsilon(n)\), and its barotropic pressure is
\(p_R=n\varepsilon'(n)-\varepsilon(n)\). Source-free homogeneous continuity
then gives \(n a^3=\mathrm{constant}\). For item (iii), spherical integration
of
\(\nabla\cdot(\kappa_R|\nabla\theta|\nabla\theta)=\beta_b\rho_b\)
gives
\(4\pi r^2\kappa_R|\theta'|\theta'=\beta_bM_b\); the stated acceleration and
definition of \(a_0\) give the result. Item (iv) follows by translating the
source term in Eq.~\eqref{eq:rpc-action} and applying translation invariance
to matter plus repair field.
\end{proof}

Define the integrated device repair charge
\begin{equation}
Q_R^{\rm dev}=q_\star\chican
\int_\Omega\scohcan(x)\dd^3x.
\label{eq:device-charge}
\end{equation}
For a prescribed external repair phase, translation of the source gives
\begin{equation}
\boxed{
F_i^{\rm coh}=-q_\star\chican
\int_\Omega\scohcan(x)\,
\partial_i\theta_{\rm ext}(x)\dd^3x.}
\label{eq:force}
\end{equation}
If the external gradient varies slowly across the device,
\begin{equation}
\boxed{F_i^{\rm dev}\simeq
-Q_R^{\rm dev}\partial_i\theta_{\rm ext}.}
\label{eq:slow-force}
\end{equation}
The sign of \(Q_R^{\rm dev}\) may select attractive or repulsive response.
Positive repair energy is compatible with signed charge when
\(\varepsilon(n)=\varepsilon(|n|)\).

\begin{remark}[Switchable-sign boundary]
\label{rem:switchable-sign}
The gated source in Eq.~\eqref{eq:source} is nonnegative under its displayed
definition. With fixed \(q_\star\) and positive \(\chican\), exchanging top and
bottom source profiles does not reverse the monopole
\(Q_R^{\rm dev}\). A switchable force sign therefore requires either a signed
orientation factor derived as part of the coherent source map or a controlled
reversal of \(\nabla\theta_{\rm ext}\). The finite channel theorem supplies
neither operation.
\end{remark}

\section{Compact-Source Constraint}

\begin{proposition}[No monopole from a neutral compact source]
Suppose the coherent repair source is compactly supported and satisfies
\begin{equation}
\int_\Omega\mathcal Q_{\rm coh}\dd^3x=0.
\label{eq:neutral}
\end{equation}
Then it has no repair monopole. Its exterior field begins at dipole order or
higher. In a uniform external gradient, the leading forces on its positive and
negative charges cancel.
\end{proposition}

\begin{proof}
Integrating Eq.~\eqref{eq:continuity} in a static state makes the exterior
flux equal the integrated source. Equation~\eqref{eq:neutral} sets the
monopole flux to zero. The multipole expansion therefore begins beyond the
monopole. A uniform external gradient couples only to total charge at leading
order, which vanishes.
\end{proof}

A bottom-to-top contrast
\begin{equation}
\Delta S=S_{\rm bottom}-S_{\rm top}
\end{equation}
can shape a dipole, select the sign of injected repair charge, couple to a
nonuniform external field, or drive an alternating repair current. It does not
by itself define a nonzero \(Q_R^{\rm dev}\). No formula proportional only to
\(g^2A\chinu\Delta S/(4\pi G)\) is a force theorem for a closed compact
device.

\section{Momentum and Energy Closure}

For the cubic condensed branch, the repair-field stress is
\begin{equation}
T^R_{ij}=\kappa_R|\nabla\theta|\,
\partial_i\theta\,\partial_j\theta
-\delta_{ij}\frac{\kappa_R}{3}|\nabla\theta|^3+\cdots .
\label{eq:stress}
\end{equation}
The exact force ledger is
\begin{equation}
\boxed{
F_i^{\rm matter}+\frac{\dd P_i^R}{\dd t}
+\oint_{\partial\Omega}T^R_{ij}n^j\dd A=0.}
\label{eq:momentum}
\end{equation}
A local repulsive response requires either nonzero device repair charge with a
field tail or explicit momentum transfer through
Eq.~\eqref{eq:momentum}. A static, closed, repair-neutral device cannot lift
itself.

Any switchable force also requires a toggle-energy ledger. Let \(E(q,s)\) be
the energy at position \(q\) and internal state \(s\), and define
\begin{equation}
\operatorname{Toggle}(q)=E(q,{\rm on})-E(q,{\rm off}).
\end{equation}
For an on/off spatial cycle between \(q_1\) and \(q_2\), conservation gives
\begin{equation}
W_{\rm cycle}=\operatorname{Toggle}(q_1)
-\operatorname{Toggle}(q_2).
\end{equation}
A claimed position-dependent switchable force with no corresponding toggle or
field-energy difference violates the closed energy ledger.

\section{Laboratory Evidence Conditions}

A force experiment has physical standing only if it supplies:
\begin{enumerate}
\item an instrument-independent measurement of \(\scohcan\) or its declared
engineering coordinate;
\item the dimensional calibration \(q_\star\);
\item a measured ambient repair gradient or a controlled emitted repair field;
\item a nonzero integrated repair charge, or a measured outgoing repair-current
and momentum flux;
\item a complete toggle, maintenance-power, and field-energy ledger;
\item detuned, dummy, reversed-sign, thermal, acoustic, electromagnetic,
mechanical, and center-of-mass controls;
\item a force result that follows the repair-field geometry and sign law in
Eq.~\eqref{eq:force}.
\item for any claimed switchable reversal, a derived signed source factor or
an independently measured reversal of the external repair gradient.
\end{enumerate}

Without these receipts, a null result does not measure \(\chican\), and a
nonzero balance signal does not establish repair charge. The lane remains
outside the falsification program because the action, source calibration, and
operational source observable are work in progress.

\section{Conclusion}

The canonical coherent-matter susceptibility is an order-one channel
coefficient fixed by protected-reserve survival on the stated OPH branch. The
repair-charge condensate completion gives that coefficient a lawful dynamical
role: it weights the coherent source in the repair-current equation. The same
proposed action yields exact finite continuity, a dilute dust limit, the
spherical deep-galaxy scaling, repair-field stress, and a device-force law.
A device force depends on the dimensional calibration, integrated repair
charge, external field, and field-momentum ledger. Repulsive response is
conditionally allowed. Reactionless lift from a compact neutral contrast is
excluded.

\begin{thebibliography}{99}

\bibitem{oph-synthesis}
B. Mueller et al.,
\emph{Observers Are All You Need}.
\url{https://github.com/FloatingPragma/observer-patch-holography/blob/main/paper/observers_are_all_you_need.pdf}

\bibitem{oph-compact}
B. Mueller et al.,
\emph{Recovering Relativity and the Standard Model from Observer Overlap Consistency}.
\url{https://github.com/FloatingPragma/observer-patch-holography/blob/main/paper/recovering_relativity_and_standard_model_structure_from_observer_overlap_consistency_compact.pdf}

\bibitem{oph-micro}
B. Mueller et al.,
\emph{Federated Echosahedral Screen Microphysics: Patch Hardware, Records, and Observer Synchronization in OPH}.
\url{https://github.com/FloatingPragma/observer-patch-holography/blob/main/paper/screen_microphysics_and_observer_synchronization.pdf}

\bibitem{oph-dark}
B. Mueller and D. Matscheko,
\emph{Observer-Patch Holography and the Dark Matter Phenomenon: A Repair-Charge Condensate Completion}.
\url{https://github.com/FloatingPragma/observer-patch-holography/blob/main/cosmology/oph_dark_matter_paper.pdf}

\end{thebibliography}

\end{document}
